% !TEX TS-program = pdflatex
\documentclass[12pt,a4paper]{article}

\usepackage[T2A]{fontenc}
\usepackage[utf8]{inputenc}
\usepackage[russian]{babel}
\usepackage{amsmath,amssymb,amsfonts,amsthm,mathtools}
\usepackage{geometry}
\usepackage{hyperref}
\usepackage{enumitem}
\usepackage{bm}

\geometry{margin=2.5cm}
\hypersetup{colorlinks=true, linkcolor=blue, citecolor=blue, urlcolor=blue}

\newtheorem{theorem}{Теорема}[section]
\newtheorem{lemma}[theorem]{Лемма}
\newtheorem{proposition}[theorem]{Предложение}
\newtheorem{corollary}[theorem]{Следствие}
\theoremstyle{definition}
\newtheorem{definition}[theorem]{Определение}
\newtheorem{remark}[theorem]{Замечание}
\newtheorem{assumption}[theorem]{Предположение}

\newcommand{\eps}{\varepsilon}
\newcommand{\R}{\mathbb{R}}
\newcommand{\N}{\mathbb{N}}
\newcommand{\Nc}{\mathbb{N}_0}
\newcommand{\Dc}{D_C^\alpha}
\newcommand{\Lc}{L_c}
\newcommand{\Rc}{R_c}
\newcommand{\Span}{\operatorname{span}}
\newcommand{\inner}[2]{\left(#1,#2\right)}
\newcommand{\norm}[1]{\left\lVert #1 \right\rVert}
\newcommand{\ealpha}[1]{E_{\alpha}\!\left(#1\right)}
\newcommand{\eab}[2]{E_{\alpha,\beta}\!\left(#1\right)} % placeholder if needed manually

\title{Резольвентно-индуцированный Müntz--Mittag--Leffler\\
спектральный Petrov--Galerkin метод\\
для сингулярно возмущённого дробного уравнения}
\author{}
\date{}

\begin{document}

\maketitle

\begin{abstract}
Предлагается новый спектральный метод для задачи
\[
\eps \Dc u(x)+a(x)u(x)=f(x), \qquad x\in(0,T], \qquad u(0)=u_0,
\]
где $0<\alpha<1$, $\Dc$ --- дробная производная Капуто, а $\eps>0$ есть малый параметр.
Ключевая идея состоит в том, чтобы:
(i) выделить главный дробный начальный слой,
(ii) перейти к скрытой плотности $h=(\eps \Dc + a_c I)(u-u_0\ell_\eps)$,
(iii) аппроксимировать $h$ в специальном спектральном пространстве, порождённом полугруппой показателей
\[
\Lambda_\alpha=\{m+n\alpha:\; m,n\in\Nc\},
\]
и
(iv) строить trial-базис в физическом пространстве как точный образ test-базиса под резольвентой модельного оператора.
Это приводит к новой операторно-согласованной спектральной схеме, допускающей
$\eps$-равномерное обоснование, квазиоптимальность и спектральную сходимость
в естественных аналитических классах.
\end{abstract}

\section{Постановка задачи}

Рассматривается задача
\begin{equation}
\label{eq:main-problem}
\eps \Dc u(x)+a(x)u(x)=f(x), \qquad x\in(0,T], \qquad u(0)=u_0,
\end{equation}
где
\[
0<\alpha<1,\qquad \eps>0,\qquad T>0.
\]
Под дробной производной Капуто понимается оператор
\begin{equation}
\label{eq:caputo}
\Dc u(x):=
\frac{1}{\Gamma(1-\alpha)}\int_0^x \frac{u'(s)}{(x-s)^\alpha}\,ds,
\qquad x\in(0,T].
\end{equation}

\begin{assumption}
\label{ass:a-f}
Предполагается, что
\[
a\in C[0,T], \qquad f\in L^2(0,T), \qquad 0<a_*\le a(x)\le a^* \quad \text{на } [0,T].
\]
\end{assumption}

Будем использовать двухпараметрическую функцию Миттаг--Леффлера
\[
E_{\alpha,\beta}(z):=\sum_{n=0}^\infty \frac{z^n}{\Gamma(\alpha n+\beta)},
\qquad \alpha>0,\ \beta>0,
\]
и обозначение $E_\alpha(z):=E_{\alpha,1}(z)$.

Основная трудность задачи \eqref{eq:main-problem} состоит в одновременном наличии
дробной не локальности и сингулярного возмущения по малому параметру $\eps$.
Даже для гладких $a$ и $f$ решение обычно имеет выраженный начальный слой
дробного типа, плохо согласованный со стандартными полиномиальными спектральными базисами.

\section{Выделение главного дробного слоя}

Положим
\[
a_0:=a(0).
\]
Определим функцию главного слоя
\begin{equation}
\label{eq:layer}
\ell_\eps(x):=E_\alpha\!\left(-\frac{a_0}{\eps}x^\alpha\right).
\end{equation}
Из стандартного тождества для производной Капуто следует, что
\begin{equation}
\label{eq:layer-eq}
\eps \Dc \ell_\eps(x)+a_0\ell_\eps(x)=0,
\qquad \ell_\eps(0)=1.
\end{equation}

Ищем решение задачи \eqref{eq:main-problem} в виде
\begin{equation}
\label{eq:u-decomp}
u(x)=u_0\ell_\eps(x)+v(x),
\qquad v(0)=0.
\end{equation}
Подставляя \eqref{eq:u-decomp} в \eqref{eq:main-problem} и используя \eqref{eq:layer-eq}, получаем
\begin{equation}
\label{eq:v-eq}
\eps \Dc v(x)+a(x)v(x)=g_\eps(x),
\qquad v(0)=0,
\end{equation}
где
\begin{equation}
\label{eq:g-eps}
g_\eps(x):=f(x)-u_0\bigl(a(x)-a_0\bigr)\ell_\eps(x).
\end{equation}

Таким образом, главный начальный слой вынесен явно, а остаток $v$ имеет существенно более
регулярную структуру.

\section{Резольвентная редукция к уравнению второго рода}

Выберем опорный коэффициент
\begin{equation}
\label{eq:ac}
a_c>0,
\qquad \text{например } \quad a_c:=\frac{a_*+a^*}{2}.
\end{equation}
Введём модельный оператор
\begin{equation}
\label{eq:Lc}
\Lc := \eps \Dc + a_c I.
\end{equation}
Определим \emph{скрытую плотность}
\begin{equation}
\label{eq:hidden-density}
h:=\Lc v = \eps \Dc v + a_c v.
\end{equation}

\begin{lemma}[Резольвента модельного оператора]
\label{lem:resolvent}
Для любого $h\in L^2(0,T)$ задача
\[
\Lc v = h, \qquad v(0)=0,
\]
имеет единственное решение
\begin{equation}
\label{eq:resolvent}
v(x)=(\Rc h)(x):=
\frac{1}{\eps}\int_0^x (x-s)^{\alpha-1}
E_{\alpha,\alpha}\!\left(-\frac{a_c}{\eps}(x-s)^\alpha\right)h(s)\,ds.
\end{equation}
\end{lemma}

\begin{proof}[План доказательства]
Применяется преобразование Лапласа по переменной $x$ к задаче с нулевым начальным условием.
В образах получаем
\[
\bigl(\eps s^\alpha+a_c\bigr)\widehat{v}(s)=\widehat{h}(s),
\qquad
\widehat{v}(s)=\frac{\widehat{h}(s)}{\eps s^\alpha+a_c}.
\]
Затем используется известная формула обратного преобразования Лапласа
\[
\mathcal{L}^{-1}\left\{\frac{1}{\eps s^\alpha+a_c}\right\}(x)
=
\frac{1}{\eps}x^{\alpha-1}E_{\alpha,\alpha}\!\left(-\frac{a_c}{\eps}x^\alpha\right),
\]
откуда следует \eqref{eq:resolvent}. Единственность вытекает из линейности и положительности
оператора $\eps \Dc + a_c I$ при $a_c>0$.
\end{proof}

Используя \eqref{eq:v-eq} и \eqref{eq:hidden-density}, получаем
\[
h+(a-a_c)v=g_\eps.
\]
Так как $v=\Rc h$, то скрытая плотность удовлетворяет интегральному уравнению второго рода
\begin{equation}
\label{eq:density-eq}
h(x)+\bigl(a(x)-a_c\bigr)(\Rc h)(x)=g_\eps(x), \qquad x\in(0,T].
\end{equation}

\begin{proposition}[Эквивалентность]
\label{prop:equivalence}
Пусть выполнено Предположение~\ref{ass:a-f}. Тогда следующие утверждения эквивалентны:
\begin{enumerate}[label=\textup{(\roman*)}]
\item $u$ решает задачу \eqref{eq:main-problem};
\item функция $h$, определённая по формулам \eqref{eq:u-decomp}, \eqref{eq:hidden-density},
решает уравнение \eqref{eq:density-eq}, а
\begin{equation}
\label{eq:u-recover}
u(x)=u_0\ell_\eps(x)+(\Rc h)(x).
\end{equation}
\end{enumerate}
\end{proposition}

\begin{proof}[План доказательства]
Нужно выполнить прямую и обратную подстановку.
В прямом направлении из \eqref{eq:main-problem} после вычитания главного слоя
получается \eqref{eq:v-eq}; затем через определение $h$ и представление $v=\Rc h$
получаем \eqref{eq:density-eq}.
В обратном направлении из \eqref{eq:density-eq} строится $v=\Rc h$,
после чего \eqref{eq:u-recover} удовлетворяет \eqref{eq:main-problem}
в силу равенств \eqref{eq:layer-eq} и \eqref{eq:hidden-density}.
\end{proof}

\section{Полугруппа показателей и новый спектральный базис}

\subsection{Мотивация выбора показателей}

Формальное локальное разложение гладкого коэффициента $a(x)$ порождает целые степени $x^m$,
а слой \eqref{eq:layer} имеет разложение
\[
\ell_\eps(x)=\sum_{n=0}^\infty
\frac{(-a_0/\eps)^n}{\Gamma(\alpha n+1)}x^{n\alpha}.
\]
Следовательно, смешанные члены вида $(a(x)-a_0)\ell_\eps(x)$ и аналогичные поправки
естественно порождают показатели
\[
m+n\alpha,\qquad m,n\in\Nc.
\]
Это мотивирует введение аддитивной полугруппы
\begin{equation}
\label{eq:Lambda}
\Lambda_\alpha:=\{m+n\alpha:\; m,n\in\Nc\}.
\end{equation}
Обозначим через
\[
0=\lambda_0<\lambda_1<\lambda_2<\dots
\]
строго возрастающую нумерацию \emph{различных} элементов множества $\Lambda_\alpha$.

\subsection{M\"untz-базис в пространстве плотности}

Рассмотрим систему функций
\[
x^{\lambda_0},\ x^{\lambda_1},\ x^{\lambda_2},\dots
\]
и ортонормируем её в $L^2(0,T)$ методом Грама--Шмидта. Получим базис
\[
M_0^{(\alpha)}, M_1^{(\alpha)},\dots,
\qquad
M_j^{(\alpha)}(x)=\sum_{k=0}^j q_{jk}x^{\lambda_k},
\]
удовлетворяющий
\[
\inner{M_i^{(\alpha)}}{M_j^{(\alpha)}}=\delta_{ij},
\qquad
\inner{p}{q}:=\int_0^T p(x)q(x)\,dx.
\]
Матрица Грама известна в явном виде:
\begin{equation}
\label{eq:Gram}
G_{ij}=\int_0^T x^{\lambda_i+\lambda_j}\,dx
=
\frac{T^{\lambda_i+\lambda_j+1}}{\lambda_i+\lambda_j+1}.
\end{equation}

Для каждого $N\in\N$ определим test-пространство
\begin{equation}
\label{eq:WN}
W_N^\alpha:=\Span\{M_0^{(\alpha)},\dots,M_{N-1}^{(\alpha)}\}.
\end{equation}

\subsection{Операторно-индуцированный trial-базис}

\begin{lemma}[Точное действие резольвенты на степенные функции]
\label{lem:monomial-image}
Для любого $\lambda>-1$ выполнено
\begin{equation}
\label{eq:monomial-image}
\Rc[x^\lambda](x)
=
\frac{\Gamma(\lambda+1)}{\eps}
x^{\lambda+\alpha}
E_{\alpha,\lambda+\alpha+1}\!\left(-\frac{a_c}{\eps}x^\alpha\right).
\end{equation}
\end{lemma}

\begin{proof}[План доказательства]
В формуле \eqref{eq:resolvent} подставляется $h(s)=s^\lambda$.
После разложения ядра в ряд
\[
E_{\alpha,\alpha}(z)=\sum_{n=0}^\infty \frac{z^n}{\Gamma(\alpha n+\alpha)}
\]
интеграл вычисляется почленно с помощью формулы Бета:
\[
\int_0^x (x-s)^{\mu-1}s^\lambda\,ds
=
x^{\mu+\lambda}B(\mu,\lambda+1)
=
x^{\mu+\lambda}\frac{\Gamma(\mu)\Gamma(\lambda+1)}{\Gamma(\mu+\lambda+1)}.
\]
Собирая ряд, получаем именно \eqref{eq:monomial-image}.
\end{proof}

Определим новые trial-функции
\begin{equation}
\label{eq:trial-basis}
\Phi_j^{\eps,\alpha}(x):=(\Rc M_j^{(\alpha)})(x), \qquad j=0,1,\dots.
\end{equation}
Из леммы \ref{lem:monomial-image} следует явная формула
\begin{equation}
\label{eq:Phi-explicit}
\Phi_j^{\eps,\alpha}(x)
=
\sum_{k=0}^j q_{jk}\,
\frac{\Gamma(\lambda_k+1)}{\eps}
x^{\lambda_k+\alpha}
E_{\alpha,\lambda_k+\alpha+1}\!\left(-\frac{a_c}{\eps}x^\alpha\right).
\end{equation}
По построению
\begin{equation}
\label{eq:intertwining}
\Lc \Phi_j^{\eps,\alpha}=M_j^{(\alpha)},\qquad j=0,1,\dots.
\end{equation}
Определим trial-пространство
\begin{equation}
\label{eq:VN}
V_N^{\eps,\alpha}:=\Rc W_N^\alpha
=
\Span\{\Phi_0^{\eps,\alpha},\dots,\Phi_{N-1}^{\eps,\alpha}\}.
\end{equation}

\begin{remark}
Соотношение \eqref{eq:intertwining} означает, что главный дробный оператор
$\Lc=\eps\Dc+a_c I$ не аппроксимируется напрямую дифференцирующей матрицей,
а \emph{точно согласован} с парой пространств
\[
V_N^{\eps,\alpha}\xrightarrow{\ \Lc\ }W_N^\alpha.
\]
Именно это отличает предлагаемую схему от стандартных полиномиальных collocation-
или Galerkin-подходов.
\end{remark}

\section{Дискретная схема}

\subsection{Форма на скрытой плотности}

Найти $h_N\in W_N^\alpha$ такое, что
\begin{equation}
\label{eq:discrete-h}
\inner{h_N + (a-a_c)\Rc h_N - g_\eps}{w_N}=0,
\qquad \forall w_N\in W_N^\alpha.
\end{equation}
После нахождения $h_N$ положим
\begin{equation}
\label{eq:uN-from-hN}
u_N(x):=u_0\ell_\eps(x)+(\Rc h_N)(x).
\end{equation}

\subsection{Эквивалентная Petrov--Galerkin форма}

Так как $h_N=\Lc(u_N-u_0\ell_\eps)$, задача \eqref{eq:discrete-h} эквивалентна:
найти
\[
u_N \in u_0\ell_\eps + V_N^{\eps,\alpha}
\]
такое, что
\begin{equation}
\label{eq:discrete-u}
\inner{\eps \Dc u_N + a u_N - f}{w_N}=0,
\qquad \forall w_N\in W_N^\alpha.
\end{equation}

\subsection{Матричная форма}

Пусть
\[
h_N(x)=\sum_{j=0}^{N-1} c_j M_j^{(\alpha)}(x),
\qquad
u_N(x)=u_0\ell_\eps(x)+\sum_{j=0}^{N-1} c_j \Phi_j^{\eps,\alpha}(x).
\]
Тогда из \eqref{eq:discrete-h} получаем систему
\begin{equation}
\label{eq:matrix-system}
A\bm{c}=\bm{b},
\end{equation}
где
\begin{equation}
\label{eq:Aij}
A_{ij}
=
\delta_{ij}
+
\inner{(a-a_c)\Phi_j^{\eps,\alpha}}{M_i^{(\alpha)}},
\qquad
b_i=\inner{g_\eps}{M_i^{(\alpha)}},
\qquad
0\le i,j\le N-1.
\end{equation}
Главная часть матрицы есть единичная матрица, а всё влияние неоднородности коэффициента
$a(x)$ сосредоточено в ограниченном возмущении.

\section{Основные теоремы}

Обозначим через $P_N$ ортогональный проектор в $L^2(0,T)$ на пространство $W_N^\alpha$.
Введём также оператор
\[
K:= (a-a_c)\Rc.
\]

\begin{theorem}[$\eps$-равномерная устойчивость непрерывной и дискретной задач]
\label{thm:stability}
Пусть выполнено Предположение~\ref{ass:a-f} и
\begin{equation}
\label{eq:rho}
\rho:=\frac{\norm{a-a_c}_{L^\infty(0,T)}}{a_c}<1.
\end{equation}
Тогда:
\begin{enumerate}[label=\textup{(\roman*)}]
\item уравнение \eqref{eq:density-eq} имеет единственное решение $h\in L^2(0,T)$;
\item для любого $N\in\N$ дискретная задача \eqref{eq:discrete-h} имеет единственное решение $h_N\in W_N^\alpha$;
\item имеют место оценки
\begin{equation}
\label{eq:stability-h}
\norm{h}_{L^2(0,T)}
\le
\frac{1}{1-\rho}\norm{g_\eps}_{L^2(0,T)},
\qquad
\norm{h_N}_{L^2(0,T)}
\le
\frac{1}{1-\rho}\norm{P_N g_\eps}_{L^2(0,T)}.
\end{equation}
\end{enumerate}
\end{theorem}

\begin{proof}[План доказательства]
Из \eqref{eq:resolvent} ядро резольвенты равно
\[
G_c^\eps(x):=\frac{1}{\eps}x^{\alpha-1}
E_{\alpha,\alpha}\!\left(-\frac{a_c}{\eps}x^\alpha\right), \qquad x>0.
\]
Используя тождество
\[
\int_0^x G_c^\eps(t)\,dt
=
\frac{1}{a_c}\left(1-E_\alpha\!\left(-\frac{a_c}{\eps}x^\alpha\right)\right)
\le \frac{1}{a_c},
\]
получаем
\[
\norm{\Rc}_{L^1\to L^1}\le \frac{1}{a_c},
\qquad
\norm{\Rc}_{L^\infty\to L^\infty}\le \frac{1}{a_c}.
\]
По интерполяции Рисса--Торина отсюда следует
\[
\norm{\Rc}_{L^2\to L^2}\le \frac{1}{a_c}.
\]
Следовательно,
\[
\norm{K}_{L^2\to L^2}
\le
\norm{a-a_c}_{L^\infty}\norm{\Rc}_{L^2\to L^2}
\le \rho<1.
\]
Тогда оператор $I+K$ обратим в $L^2(0,T)$ по ряду Неймана, что даёт существование,
единственность и первую оценку в \eqref{eq:stability-h}. Для дискретной схемы используется тот же аргумент
для оператора $I+P_NK$ на конечномерном пространстве $W_N^\alpha$.
\end{proof}

\begin{theorem}[Квазиоптимальность]
\label{thm:quasioptimal}
Пусть выполнены условия Теоремы~\ref{thm:stability}. Тогда для решения $h$ уравнения
\eqref{eq:density-eq} и решения $h_N$ задачи \eqref{eq:discrete-h} верно
\begin{equation}
\label{eq:quasi-h}
\norm{h-h_N}_{L^2(0,T)}
\le
\frac{1}{1-\rho}
\inf_{q_N\in W_N^\alpha}\norm{h-q_N}_{L^2(0,T)}.
\end{equation}
Для ошибки в исходной неизвестной имеем
\begin{equation}
\label{eq:quasi-u}
\norm{u-u_N}_{L^2(0,T)}
\le
\frac{1}{a_c(1-\rho)}
\inf_{q_N\in W_N^\alpha}\norm{h-q_N}_{L^2(0,T)}.
\end{equation}
\end{theorem}

\begin{proof}[План доказательства]
Поскольку $h$ и $h_N$ удовлетворяют соответственно
\[
h+Kh=g_\eps,
\qquad
h_N+P_NKh_N=P_Ng_\eps,
\]
то для ошибки проекции $P_Nh-h_N$ получаем
\[
(I+P_NK)(P_Nh-h_N)=-P_NK(h-P_Nh).
\]
Из обратимости $I+P_NK$ и оценки $\norm{P_NK}\le \rho$ следует
\[
\norm{P_Nh-h_N}\le \frac{\rho}{1-\rho}\norm{h-P_Nh}.
\]
Далее
\[
\norm{h-h_N}\le \norm{h-P_Nh}+\norm{P_Nh-h_N}
\le \frac{1}{1-\rho}\norm{h-P_Nh},
\]
что и даёт \eqref{eq:quasi-h}, так как $P_Nh$ есть наилучшее приближение в $W_N^\alpha$.
Наконец,
\[
u-u_N=\Rc(h-h_N),
\]
и оценка \eqref{eq:quasi-u} следует из неравенства $\norm{\Rc}_{L^2\to L^2}\le 1/a_c$.
\end{proof}

\begin{definition}[M\"untz-аналитический класс]
\label{def:analytic-class}
Для $\sigma>0$ определим пространство
\[
\mathcal{A}_{\Lambda_\alpha,\sigma}
:=
\left\{
z=\sum_{j=0}^\infty \widehat z_j M_j^{(\alpha)} \;:\;
\sum_{j=0}^\infty e^{2\sigma \lambda_j}|\widehat z_j|^2<\infty
\right\},
\]
с нормой
\[
\norm{z}_{\mathcal{A}_{\Lambda_\alpha,\sigma}}^2
:=
\sum_{j=0}^\infty e^{2\sigma \lambda_j}|\widehat z_j|^2.
\]
\end{definition}

\begin{theorem}[Спектральная сходимость в M\"untz-аналитическом классе]
\label{thm:spectral}
Пусть выполнены условия Теоремы~\ref{thm:stability} и скрытая плотность
$h$ принадлежит классу $\mathcal{A}_{\Lambda_\alpha,\sigma}$ для некоторого $\sigma>0$.
Тогда
\begin{equation}
\label{eq:spectral-h}
\norm{h-h_N}_{L^2(0,T)}
\le
\frac{e^{-\sigma \lambda_N}}{1-\rho}
\norm{h}_{\mathcal{A}_{\Lambda_\alpha,\sigma}},
\end{equation}
а для исходной неизвестной
\begin{equation}
\label{eq:spectral-u}
\norm{u-u_N}_{L^2(0,T)}
\le
\frac{e^{-\sigma \lambda_N}}{a_c(1-\rho)}
\norm{h}_{\mathcal{A}_{\Lambda_\alpha,\sigma}}.
\end{equation}
\end{theorem}

\begin{proof}[План доказательства]
Поскольку $\{M_j^{(\alpha)}\}_{j\ge0}$ --- ортонормированная система, для ортогонального проектора $P_N$
имеем
\[
\norm{h-P_Nh}_{L^2(0,T)}^2
=
\sum_{j=N}^\infty |\widehat h_j|^2
\le
e^{-2\sigma \lambda_N}
\sum_{j=N}^\infty e^{2\sigma \lambda_j}|\widehat h_j|^2
\le
e^{-2\sigma \lambda_N}
\norm{h}_{\mathcal{A}_{\Lambda_\alpha,\sigma}}^2.
\]
Подстановка этой оценки в \eqref{eq:quasi-h} и \eqref{eq:quasi-u} завершает доказательство.
\end{proof}

\begin{theorem}[Асимптотически согласованный предел при $\eps\to0$]
\label{thm:AP}
Пусть выполнены условия Теоремы~\ref{thm:stability} и $N\in\N$ фиксировано.
Пусть $u_N^\eps$ есть решение дискретной схемы \eqref{eq:discrete-u}. Тогда при $\eps\to0$
\[
u_N^\eps \to \overline u_N \quad \text{в } L^2(0,T),
\]
где $\overline u_N\in W_N^\alpha$ определяется как решение предельной спектральной задачи
\begin{equation}
\label{eq:limit-PG}
\inner{a\,\overline u_N - f}{w_N}=0,
\qquad \forall w_N\in W_N^\alpha.
\end{equation}
Следовательно, схема является асимптотически согласованной
(\emph{asymptotic-preserving}) для вырожденного предела $\eps\to0$.
\end{theorem}

\begin{proof}[План доказательства]
Сначала показывается, что
\[
\ell_\eps \to 0 \quad \text{в } L^2(0,T),
\]
поскольку дробный начальный слой концентрируется в сужающейся окрестности точки $x=0$.
Затем из формулы \eqref{eq:resolvent} или из образов Лапласа следует сильная сходимость
\[
\Rc \to a_c^{-1}I \qquad \text{в } \mathcal{L}(L^2(0,T),L^2(0,T)).
\]
В частности,
\[
\Phi_j^{\eps,\alpha}=\Rc M_j^{(\alpha)} \to a_c^{-1}M_j^{(\alpha)}
\qquad \text{в } L^2(0,T)
\]
для каждого фиксированного $j$.
Кроме того, из \eqref{eq:g-eps} имеем
\[
g_\eps \to f \qquad \text{в } L^2(0,T).
\]
Поэтому матрица дискретной системы \eqref{eq:matrix-system} и правая часть сходятся:
\[
A^\eps_{ij}\to \delta_{ij}+\inner{(a-a_c)a_c^{-1}M_j^{(\alpha)}}{M_i^{(\alpha)}},
\qquad
b_i^\eps\to \inner{f}{M_i^{(\alpha)}}.
\]
Предельная система эквивалентна \eqref{eq:limit-PG}. Конечномерность пространства $W_N^\alpha$
обеспечивает сходимость коэффициентов и, следовательно, сходимость $u_N^\eps$ в $L^2(0,T)$.
\end{proof}

\begin{proposition}[Точность на модельной задаче с постоянным коэффициентом]
\label{prop:exact-constant}
Если $a(x)\equiv a_c$ на $[0,T]$ и $g_\eps\in W_N^\alpha$, то дискретное решение совпадает
с точным:
\[
h_N=h=g_\eps,
\qquad
u_N=u.
\]
\end{proposition}

\begin{proof}[План доказательства]
При $a\equiv a_c$ уравнение \eqref{eq:density-eq} упрощается до $h=g_\eps$.
Если $g_\eps\in W_N^\alpha$, то из \eqref{eq:discrete-h} следует $h_N=P_Ng_\eps=g_\eps$.
Далее $u_N=u_0\ell_\eps+\Rc h_N=u_0\ell_\eps+\Rc h=u$.
\end{proof}

\section{Итоговая формулировка метода}

Предлагаемый метод можно кратко сформулировать так:

\begin{enumerate}[label=\textbf{Шаг \arabic*.}]
\item Выделить главный слой
\[
u=u_0\ell_\eps+v,
\qquad
\ell_\eps(x)=E_\alpha\!\left(-\frac{a_0}{\eps}x^\alpha\right).
\]

\item Ввести скрытую плотность
\[
h=(\eps \Dc + a_c I)v
\]
и свести задачу к уравнению второго рода
\[
h+(a-a_c)\Rc h=g_\eps.
\]

\item Построить M\"untz-test-пространство $W_N^\alpha$ по показателям
\[
\Lambda_\alpha=\Nc+\alpha\Nc
\]
и ортонормированным функциям $M_j^{(\alpha)}$.

\item Построить operator-induced trial-пространство
\[
V_N^{\eps,\alpha}=\Rc W_N^\alpha
=
\Span\{\Phi_j^{\eps,\alpha}\}_{j=0}^{N-1},
\]
где
\[
\Phi_j^{\eps,\alpha}=\Rc M_j^{(\alpha)}
\]
заданы явно формулой \eqref{eq:Phi-explicit}.

\item Решить Petrov--Galerkin систему
\[
\inner{\eps \Dc u_N + a u_N - f}{w_N}=0,
\qquad \forall w_N\in W_N^\alpha.
\]
\end{enumerate}

Концептуальная новизна конструкции состоит в том, что trial-базис
не выбирается из заранее известных семейств полиномов или рациональных функций,
а порождается \emph{самим модельным дробным оператором} через резольвенту.
Это делает схему одновременно
\begin{itemize}
\item согласованной с дробной структурой,
\item приспособленной к сингулярному начальному слою,
\item удобной для $\eps$-равномерного анализа устойчивости,
\item пригодной для спектральной сходимости в естественных M\"untz-аналитических классах.
\end{itemize}

\section{Модельные задачи для численной проверки метода}
\label{sec:tests}

В этом разделе предлагается набор задач, позволяющих проверить различные аспекты
предлагаемого резольвентно-индуцированного Müntz--Mittag--Leffler метода:
\begin{enumerate}[label=\textup{(\roman*)}]
\item алгебраическую корректность построения trial/test-пространств;
\item спектральную сходимость на решениях, согласованных с полугруппой показателей
$\Lambda_\alpha=\Nc+\alpha\Nc$;
\item устойчивость при $\eps\ll 1$;
\item поведение на решениях, не принадлежащих естественному Müntz-классу;
\item асимптотически согласованное поведение при $\eps\to0$.
\end{enumerate}

Во всех тестах удобно, если не оговорено иное, брать
\[
T=1,\qquad
\alpha\in\{0.25,\,0.5,\,0.75\},\qquad
\eps\in\{10^{-1},\,10^{-3},\,10^{-5},\,10^{-7}\},
\]
а размерность спектрального пространства варьировать, например, по сетке
\[
N=4,6,8,10,12,\dots.
\]

\subsection{Универсальная manufactured-конструкция}

Для построения тестов с точным решением удобно использовать две общие схемы.

\paragraph{(A) Manufactured-решения степенного типа.}
Пусть
\begin{equation}
\label{eq:test-manufactured-power}
u_{\mathrm{ex}}(x)
=
u_0\ell_\eps(x)+\sum_{r=1}^R c_r x^{\mu_r},
\qquad
\mu_r>0,
\end{equation}
где
\[
\ell_\eps(x)=E_\alpha\!\left(-\frac{a_0}{\eps}x^\alpha\right),
\qquad
a_0=a(0).
\]
Тогда, используя формулу
\[
\Dc x^\mu
=
\frac{\Gamma(\mu+1)}{\Gamma(\mu+1-\alpha)}x^{\mu-\alpha},
\qquad \mu>0,
\]
получаем правую часть
\begin{equation}
\label{eq:test-manufactured-force}
f(x)
=
u_0\bigl(a(x)-a_0\bigr)\ell_\eps(x)
+
\sum_{r=1}^R c_r
\left[
\eps \frac{\Gamma(\mu_r+1)}{\Gamma(\mu_r+1-\alpha)}x^{\mu_r-\alpha}
+
a(x)x^{\mu_r}
\right].
\end{equation}
Тогда $u_{\mathrm{ex}}$ является точным решением исходной задачи.

\paragraph{(B) Manufactured-решения в operator-induced trial-пространстве.}
Пусть
\begin{equation}
\label{eq:test-manufactured-trial}
u_{\mathrm{ex}}(x)
=
u_0\ell_\eps(x)+\sum_{j=0}^{M} d_j \Phi_j^{\eps,\alpha}(x),
\end{equation}
где $\Phi_j^{\eps,\alpha}=\Rc M_j^{(\alpha)}$.
Тогда в силу тождества
\[
(\eps \Dc + a I)\Phi_j^{\eps,\alpha}
=
M_j^{(\alpha)} + (a-a_c)\Phi_j^{\eps,\alpha}
\]
получаем
\begin{equation}
\label{eq:test-force-trial}
f(x)
=
u_0\bigl(a(x)-a_0\bigr)\ell_\eps(x)
+
\sum_{j=0}^{M} d_j
\left(
M_j^{(\alpha)}(x)
+
\bigl(a(x)-a_c\bigr)\Phi_j^{\eps,\alpha}(x)
\right).
\end{equation}
Такой тест особенно важен, поскольку проверяет корректность именно нового
operator-induced базиса.

\subsection{Задача 1: точное воспроизведение operator-induced trial-пространства}

Рассмотрим
\[
a(x)=2+x,
\qquad
a_0=2,
\qquad
a_c=\frac{2+3}{2}=\frac52.
\]
Зададим точное решение в виде
\begin{equation}
\label{eq:test1-uex}
u_{\mathrm{ex}}(x)
=
u_0\ell_\eps(x)
+
\Phi_0^{\eps,\alpha}(x)
-2\Phi_2^{\eps,\alpha}(x)
+\frac12 \Phi_4^{\eps,\alpha}(x),
\end{equation}
а правую часть определим по формуле \eqref{eq:test-force-trial}, то есть
\begin{equation}
\label{eq:test1-f}
\begin{aligned}
f(x)
&=
u_0\bigl(a(x)-2\bigr)\ell_\eps(x)
+
M_0^{(\alpha)}(x)-2M_2^{(\alpha)}(x)+\frac12 M_4^{(\alpha)}(x)
\\
&\quad
+
\left(a(x)-\frac52\right)
\left(
\Phi_0^{\eps,\alpha}(x)
-2\Phi_2^{\eps,\alpha}(x)
+\frac12 \Phi_4^{\eps,\alpha}(x)
\right).
\end{aligned}
\end{equation}

\medskip
\noindent
\textbf{Что проверяет эта задача.}
Если используется точная или достаточно точная квадратурная сборка матрицы,
то при $N\ge 5$ численное решение должно воспроизводить \eqref{eq:test1-uex}
с точностью до машинного округления. Это является базовым тестом
внутренней алгебраической согласованности метода.

\subsection{Задача 2: \texorpdfstring{$\Lambda_\alpha$}{Lambda}-согласованное точное решение}

Рассмотрим коэффициент
\[
a(x)=1+x+x^2,
\qquad
a_0=1.
\]
Положим
\begin{equation}
\label{eq:test2-uex}
u_{\mathrm{ex}}(x)
=
\ell_\eps(x)
+
x^\alpha
-\frac12 x
+\frac14 x^{1+\alpha}
-\frac13 x^{2\alpha}.
\end{equation}
Все показатели
\[
\alpha,\quad 1,\quad 1+\alpha,\quad 2\alpha
\]
принадлежат полугруппе $\Lambda_\alpha=\Nc+\alpha\Nc$.
Следовательно, это естественная модельная задача для предлагаемого метода.

Правая часть строится по формуле \eqref{eq:test-manufactured-force}
с параметрами
\[
(\mu_1,c_1)=(\alpha,1),\quad
(\mu_2,c_2)=(1,-1/2),\quad
(\mu_3,c_3)=(1+\alpha,1/4),\quad
(\mu_4,c_4)=(2\alpha,-1/3).
\]
Явно:
\begin{equation}
\label{eq:test2-f}
\begin{aligned}
f(x)
&=
(x+x^2)\ell_\eps(x)
+
\eps\Gamma(\alpha+1)
-\frac{\eps}{2}\frac{1}{\Gamma(2-\alpha)}x^{1-\alpha}
\\
&\quad
+\frac{\eps}{4}\frac{\Gamma(2+\alpha)}{\Gamma(2)}x
-\frac{\eps}{3}\frac{\Gamma(2\alpha+1)}{\Gamma(\alpha+1)}x^\alpha
\\
&\quad
+(1+x+x^2)
\left(
x^\alpha
-\frac12 x
+\frac14 x^{1+\alpha}
-\frac13 x^{2\alpha}
\right).
\end{aligned}
\end{equation}

\medskip
\noindent
\textbf{Что проверяет эта задача.}
На этой задаче следует ожидать спектральной или, по крайней мере,
супералгебраической сходимости по $N$, поскольку решение согласовано
с естественным набором показателей метода.

\subsection{Задача 3: решение вне естественного Müntz-класса}

Для проверки робастности метода на ``несогласованных'' решениях рассмотрим
\[
a(x)=1+x,
\qquad
a_0=1,
\]
и зададим точное решение
\begin{equation}
\label{eq:test3-uex}
u_{\mathrm{ex}}(x)
=
\ell_\eps(x)+x^\beta,
\qquad
\beta=\sqrt{2}.
\end{equation}
Тогда
\begin{equation}
\label{eq:test3-f}
f(x)
=
x\ell_\eps(x)
+
\eps \frac{\Gamma(\beta+1)}{\Gamma(\beta+1-\alpha)}x^{\beta-\alpha}
+
(1+x)x^\beta.
\end{equation}

\medskip
\noindent
\textbf{Что проверяет эта задача.}
Показатель $\beta=\sqrt2$ не согласован с полугруппой $\Lambda_\alpha$,
поэтому на этой задаче следует ожидать более медленной сходимости.
Тем не менее метод должен сохранять устойчивость и демонстрировать
предсказуемое ухудшение скорости, а не потерю работоспособности.

\subsection{Задача 4: \texorpdfstring{$\eps$}{eps}-равномерное разрешение начального слоя}

Для проверки равномерности по малому параметру $\eps$ рекомендуется использовать задачу~2
при серии значений
\[
\eps=10^{-1},\ 10^{-3},\ 10^{-5},\ 10^{-7}.
\]
При этом следует измерять ошибки
\[
E_N^{(2)}(\eps):=\norm{u_{\mathrm{ex}}-u_N}_{L^2(0,1)},
\qquad
E_N^{(\infty)}(\eps):=\norm{u_{\mathrm{ex}}-u_N}_{L^\infty(0,1)},
\]
для фиксированных $N$ и анализировать, насколько эти ошибки устойчивы
при уменьшении $\eps$.

Кроме того, полезно визуализировать решение в растянутой переменной
\[
\xi=\frac{x}{\eps^{1/\alpha}},
\]
поскольку именно на масштабе $x=O(\eps^{1/\alpha})$ формируется дробный начальный слой.

\medskip
\noindent
\textbf{Что проверяет эта задача.}
Данный тест проверяет, не деградирует ли аппроксимация при сжатии слоя
в окрестность точки $x=0$, и подтверждает ли схема заявленную
$\eps$-равномерную устойчивость.

\subsection{Задача 5: проверка asymptotic-preserving свойства}

Возьмём ту же постановку, что и в задаче~2. Тогда при $\eps\to0$ точное решение
\eqref{eq:test2-uex} стремится в $L^2(0,1)$ к предельной функции
\begin{equation}
\label{eq:test5-u0}
u^{0}(x)
=
x^\alpha
-\frac12 x
+\frac14 x^{1+\alpha}
-\frac13 x^{2\alpha},
\end{equation}
а предельная правая часть равна
\begin{equation}
\label{eq:test5-f0}
f^{0}(x)=a(x)u^0(x).
\end{equation}

Для фиксированного $N$ следует решать предельную дискретную задачу:
найти $\overline u_N\in W_N^\alpha$ такую, что
\begin{equation}
\label{eq:test5-limit-discrete}
\inner{a\overline u_N-f^0}{w_N}=0,
\qquad \forall w_N\in W_N^\alpha,
\end{equation}
и затем сравнивать её с решением полной задачи:
\begin{equation}
\label{eq:test5-AP-error}
A_N(\eps):=\norm{u_N^\eps-\overline u_N}_{L^2(0,1)}.
\end{equation}

\medskip
\noindent
\textbf{Что проверяет эта задача.}
Если метод действительно является asymptotic-preserving, то при фиксированном $N$
должно наблюдаться
\[
A_N(\eps)\to 0 \qquad \text{при } \eps\to0.
\]

\subsection{Задача 6: зависимость от величины возмущения коэффициента}

Для проверки роли параметра
\[
\rho=\frac{\norm{a-a_c}_{L^\infty}}{a_c}
\]
в теории устойчивости полезно рассмотреть семейство коэффициентов
\[
a_\delta(x)=1+\delta(2x-1),
\qquad 0<\delta<1.
\]
Здесь
\[
a_0=1-\delta,
\qquad
a_c=1,
\qquad
\rho=\delta.
\]
В качестве точного решения можно взять
\begin{equation}
\label{eq:test6-uex}
u_{\mathrm{ex}}(x)=\ell_\eps(x)+x,
\qquad
\ell_\eps(x)=E_\alpha\!\left(-\frac{1-\delta}{\eps}x^\alpha\right),
\end{equation}
а правая часть тогда равна
\begin{equation}
\label{eq:test6-f}
f(x)
=
2\delta x\,\ell_\eps(x)
+
\eps \frac{1}{\Gamma(2-\alpha)}x^{1-\alpha}
+
a_\delta(x)x.
\end{equation}

Следует провести вычисления, например, для
\[
\delta\in\{0.2,\,0.5,\,0.8,\,0.95\}.
\]

\medskip
\noindent
\textbf{Что проверяет эта задача.}
Данный тест позволяет эмпирически изучить, как с ростом $\rho$ меняются:
\begin{enumerate}[label=\textup{(\alph*)}]
\item точность;
\item число обусловленности матрицы $A$;
\item устойчивость вычислений;
\item скорость сходимости итерационных решателей (если они используются).
\end{enumerate}

\subsection{Рекомендуемые численные характеристики}

Для задач с известным точным решением рекомендуется вычислять:
\begin{align}
E_N^{(2)} &:= \norm{u_{\mathrm{ex}}-u_N}_{L^2(0,T)}, \\
E_N^{(\infty)} &:= \norm{u_{\mathrm{ex}}-u_N}_{L^\infty(0,T)}, \\
R_N^{(2)} &:= \norm{\eps \Dc u_N + a u_N - f}_{L^2(0,T)}.
\end{align}
Для задачи 1 дополнительно следует контролировать ошибку воспроизведения:
\[
E_{N,\mathrm{repr}}:=\norm{u_{\mathrm{ex}}-u_N}_{L^\infty(0,T)}.
\]
Для задачи 5 естественно измерять AP-ошибку \eqref{eq:test5-AP-error}.
Для задачи 6 полезно также регистрировать число обусловленности
\[
\kappa_2(A)
\]
дискретной матрицы системы.

\subsection{Ожидаемые качественные результаты}

На основании построенной теории ожидаются следующие результаты:
\begin{enumerate}[label=\textup{(\roman*)}]
\item в задаче~1 решение воспроизводится с точностью до ошибок квадратуры и машинного округления;
\item в задаче~2 наблюдается спектральное убывание ошибки по $N$;
\item в задаче~3 скорость сходимости ниже, но устойчивость сохраняется;
\item в задаче~4 при фиксированном $N$ ошибка не должна катастрофически расти при $\eps\to0$;
\item в задаче~5 должно выполняться $A_N(\eps)\to0$ при $\eps\to0$;
\item в задаче~6 ухудшение численных характеристик при росте $\rho$ должно быть плавным
и согласованным с теоретическим условием $\rho<1$.
\end{enumerate}

Тем самым предложенный набор тестов позволяет отдельно и содержательно проверить
все ключевые свойства нового метода:
точность operator-induced базиса, спектральную эффективность, устойчивость
в сингулярно возмущённом режиме и асимптотическую согласованность.

\section{Реальные прикладные задачи для валидации метода}
\label{sec:real-benchmarks}

В этом разделе вместо искусственно сконструированных \emph{manufactured solutions}
предлагаются реальные прикладные задачи, для которых в литературе уже существуют
либо экспериментальные данные, либо опубликованные параметрические идентификации,
либо устоявшиеся аналитические формулы для сравнения.
Для первой публикации по новому методу я бы рекомендовал обязательно включить
задачу~\ref{app:supercap}: она практически буквально совпадает с рассматриваемой
моделью, допускает сравнение как с точным решением через функцию Миттаг--Леффлера,
так и с опубликованными экспериментальными кривыми \cite{gomezaguilar2014rc,lopez2022cpe,kothari2021supercap,wang2020ultracap}.

\subsection{Разряд/заряд суперконденсатора и дробный RC--CPE элемент}
\label{app:supercap}

Одна из наиболее естественных прикладных постановок возникает в моделях
суперконденсаторов, батарей и RC-цепей с \emph{constant phase element} (CPE).
В простейшей одноветвевой time-domain модели напряжение $V_c(t)$ на дробно-ёмкостном
элементе удовлетворяет уравнению
\begin{equation}
\label{eq:app-supercap}
R_s C_\alpha D_t^\alpha V_c(t) + V_c(t) = V_{\mathrm{in}}(t),
\qquad
t\in(0,T],
\qquad
V_c(0)=V_0,
\end{equation}
где $R_s$ --- последовательное сопротивление, $C_\alpha$ --- дробная ёмкость,
$0<\alpha<1$, а $V_{\mathrm{in}}(t)$ --- приложенное напряжение
\cite{gomezaguilar2014rc,lopez2022cpe,kothari2021supercap}.

Эта задача \emph{в точности} приводится к исследуемой модели:
\[
u(t)=V_c(t),
\qquad
\eps = R_s C_\alpha,
\qquad
a(t)\equiv 1,
\qquad
f(t)=V_{\mathrm{in}}(t).
\]

\paragraph{Чистый разряд.}
Если после предварительного заряда источник отключён, то $V_{\mathrm{in}}(t)\equiv0$,
и получаем
\begin{equation}
\label{eq:app-supercap-discharge}
\eps D_t^\alpha u(t)+u(t)=0,
\qquad
u(0)=V_0,
\end{equation}
с точным решением
\begin{equation}
\label{eq:app-supercap-discharge-exact}
u(t)=V_0\,E_\alpha\!\left(-\frac{t^\alpha}{\eps}\right)
=
V_0\,E_\alpha\!\left(-\frac{t^\alpha}{R_s C_\alpha}\right).
\end{equation}

\paragraph{Ступенчатый заряд.}
Если на вход подаётся постоянное напряжение $V_\star$, то
\begin{equation}
\label{eq:app-supercap-charge}
\eps D_t^\alpha u(t)+u(t)=V_\star,
\qquad
u(0)=V_0,
\end{equation}
и
\begin{equation}
\label{eq:app-supercap-charge-exact}
u(t)=V_\star + \bigl(V_0-V_\star\bigr)
E_\alpha\!\left(-\frac{t^\alpha}{\eps}\right).
\end{equation}

\paragraph{Почему это хороший benchmark.}
\begin{enumerate}[label=\textup{(\roman*)}]
\item Физическая постановка совпадает с математической почти без редукции.
\item Есть точное решение через функцию Миттаг--Леффлера, поэтому можно отдельно
проверять \emph{численную корректность}.
\item Есть опубликованные time-domain и EIS-сопоставления для fractional RC/CPE
и суперконденсаторов, поэтому можно отдельно проверять \emph{прикладную релевантность}
\cite{lopez2022cpe,kothari2021supercap,wang2020ultracap}.
\item При малом $R_s C_\alpha$ и фиксированном горизонте наблюдения $T$
задача действительно работает в сингулярно возмущённом режиме.
\end{enumerate}

\paragraph{Как сравнивать с известными результатами.}
\begin{enumerate}[label=\textup{(\alph*)}]
\item Взять из литературы опубликованные параметры $(R_s,C_\alpha,\alpha)$
для конкретного суперконденсатора или CPE-модели \cite{kothari2021supercap,wang2020ultracap}.
\item Посчитать $u_N(t)$ новым методом.
\item Сравнить:
\[
\max_{t\in[0,T]} |u_N(t)-u_{\mathrm{exp}}(t)|,
\qquad
\mathrm{RMSE},
\qquad
\norm{u_N-u_{\mathrm{ML}}}_{L^\infty(0,T)},
\]
где $u_{\mathrm{exp}}$ --- экспериментальная кривая, а $u_{\mathrm{ML}}$ ---
точная кривая \eqref{eq:app-supercap-discharge-exact} или
\eqref{eq:app-supercap-charge-exact}.
\item Отдельно измерить чувствительность ошибки к уменьшению $\eps=R_sC_\alpha$.
\end{enumerate}

\subsection{Диэлектрическая релаксация биоткани или полимера после ступенчатого поля}
\label{app:dielectric}

В моделях диэлектрической релаксации неизвестной может быть поляризация $P(t)$,
которая после приложения ступенчатого электрического поля $E_0$ удовлетворяет
уравнению
\begin{equation}
\label{eq:app-dielectric}
\tau^\alpha D_t^\alpha P(t)+P(t)=\chi E_0,
\qquad
t\in(0,T],
\qquad
P(0)=P_0.
\end{equation}
Здесь $\tau$ --- характерное время релаксации, $\chi$ --- коэффициент отклика.
Именно такого типа Caputo-модели используются в работах по дробной модели
диэлектрической релаксации и сопоставляются с данными по биотканям
(в частности, human aorta) и полимерным системам \cite{ciancio2017dielectric,renteria2025polymer}.

Эта задача опять же имеет вид
\[
u(t)=P(t),
\qquad
\eps=\tau^\alpha,
\qquad
a(t)\equiv1,
\qquad
f(t)=\chi E_0.
\]
Для неё известно точное решение
\begin{equation}
\label{eq:app-dielectric-exact}
P(t)=\chi E_0 + \bigl(P_0-\chi E_0\bigr)
E_\alpha\!\left(-\frac{t^\alpha}{\tau^\alpha}\right).
\end{equation}

\paragraph{Продвинутый вариант с переменным коэффициентом.}
Если эксперимент проводится при линейном температурном прогоне
\[
T(t)=T_0+\beta t,
\]
а время релаксации зависит от температуры, например по закону Аррениуса
или по кооперативной модели, то после нормировки можно получить постановку
\begin{equation}
\label{eq:app-dielectric-variable}
\eps D_t^\alpha P(t)+a(t)P(t)=f(t),
\end{equation}
где
\[
a(t)=\frac{\eps}{\tau(T(t))^\alpha},
\qquad
f(t)=a(t)\chi E_0.
\]
Это уже реальная прикладная задача с \emph{непостоянным} коэффициентом $a(t)$,
что особенно интересно для тестирования предлагаемого метода
\cite{renteria2025polymer}.

\paragraph{Как сравнивать с известными результатами.}
\begin{enumerate}[label=\textup{(\alph*)}]
\item Взять опубликованные параметры дробной диэлектрической модели
для biological tissues или polymer datasets \cite{ciancio2017dielectric,renteria2025polymer}.
\item Решить time-domain задачу \eqref{eq:app-dielectric}
или \eqref{eq:app-dielectric-variable}.
\item Если доступны time-domain данные поляризации, сравнить с ними напрямую.
\item Если в литературе приведены frequency-domain данные по комплексной
диэлектрической проницаемости, восстановить отклик из time-domain решения, например,
через преобразование
\begin{equation}
\label{eq:app-dielectric-fourier}
\varepsilon^\ast(\omega)
=
\varepsilon_\infty
+
\frac{1}{E_0}
\int_0^\infty e^{-i\omega t}\,\dot P(t)\,dt,
\end{equation}
и сравнить положения пиков, ширины пиков и значения потерь
с опубликованными спектрами.
\end{enumerate}

\paragraph{Что именно проверяет эта задача.}
Эта постановка полезна не только как test of accuracy, но и как test of interpretation:
метод должен корректно воспроизводить очень медленную релаксацию,
начальный слой и переход к частотным характеристикам.

\subsection{Ползучесть frozen soil в fractional Kelvin--Voigt модели}
\label{app:frozen-soil}

В геомеханике и реологии для описания ползучести при постоянной нагрузке
широко используется fractional Kelvin--Voigt модель.
Если на образец действует постоянное напряжение $\sigma_0$, то деформация
$\varepsilon_m(t)$ удовлетворяет
\begin{equation}
\label{eq:app-creep}
\eta_\alpha D_t^\alpha \varepsilon_m(t)+E\,\varepsilon_m(t)=\sigma_0,
\qquad
t\in(0,T],
\qquad
\varepsilon_m(0)=0,
\end{equation}
где $E$ --- упругий модуль, а $\eta_\alpha$ --- дробный вязкоупругий параметр.
Именно такие fractional Kelvin--Voigt постановки используются для описания
ползучести frozen soil и сравниваются с результатами SIT-, uniaxial- и triaxial-creep
испытаний \cite{zhang2022frozen}.

Отождествление с нашей задачей здесь такое:
\[
u(t)=\varepsilon_m(t),
\qquad
\eps=\eta_\alpha,
\qquad
a(t)\equiv E,
\qquad
f(t)=\sigma_0.
\]
Точное решение имеет вид
\begin{equation}
\label{eq:app-creep-exact}
\varepsilon_m(t)
=
\frac{\sigma_0}{E}
\left[
1-
E_\alpha\!\left(-\frac{E}{\eta_\alpha}t^\alpha\right)
\right].
\end{equation}

\paragraph{Практический смысл этого benchmark-а.}
\begin{enumerate}[label=\textup{(\roman*)}]
\item Это настоящая time-domain механическая задача с экспериментальными кривыми
``деформация--время''.
\item Задача имеет правильное bounded-creep поведение:
\[
\varepsilon_m(t)\to \frac{\sigma_0}{E}, \qquad t\to\infty,
\]
что позволяет проверять не только начальный слой, но и правильную long-time asymptotics.
\item При малом $\eta_\alpha$ возникает быстрый начальный переход, то есть режим,
где новый метод должен проявить преимущество перед обычными time-marching схемами.
\end{enumerate}

\paragraph{Как сравнивать с известными результатами.}
\begin{enumerate}[label=\textup{(\alph*)}]
\item Взять из опубликованных frozen-soil тестов параметры $(E,\eta_\alpha,\alpha)$
или сами деформационные кривые \cite{zhang2022frozen}.
\item Для постоянного напряжения $\sigma_0$ вычислить $u_N(t)$.
\item Сравнить:
\[
\mathrm{RMSE},
\qquad
\norm{u_N-u_{\mathrm{exp}}}_{L^\infty(0,T)},
\qquad
\left|u_N(T)-\frac{\sigma_0}{E}\right|.
\]
\item Отдельно исследовать, насколько устойчиво работает метод при уменьшении
$\eta_\alpha$ и при увеличении длительности наблюдения $T$.
\end{enumerate}

\subsection{Какую задачу ставить в первую численную статью}

Если нужен \emph{один главный} прикладной benchmark, то я бы поставил на первое место
задачу~\ref{app:supercap}. Причины следующие:
\begin{enumerate}[label=\textup{(\roman*)}]
\item математическая модель почти буквально совпадает с целевым уравнением;
\item есть и точное решение, и экспериментальные кривые;
\item вопрос корректной инициализации для Caputo/CPE-моделей сам по себе нетривиален,
поэтому это хороший стресс-тест для метода;
\item в литературе fractional-order модели суперконденсаторов считаются более точными,
чем простые integer-order RC-модели, так что сравнение будет содержательным
не только математически, но и инженерно \cite{lopez2022cpe,wang2020ultracap}.
\end{enumerate}

Задачу~\ref{app:frozen-soil} я бы поставил второй: она хорошо показывает,
что метод работает не только на электротехнических примерах, но и на реологических
данных реального материала.
Задачу~\ref{app:dielectric} удобно использовать как третий пример,
особенно если хочется продемонстрировать версию с переменным коэффициентом $a(t)$
через температурный прогон.

\bibliographystyle{plain} % или ieeetr, unsrt, alpha и т.д.
\bibliography{references}

\end{document}